\documentclass[11pt]{article}
\usepackage[utf8]{inputenc}
\usepackage[a4paper,margin=1in]{geometry}
\usepackage{amsmath,amssymb}
\usepackage{booktabs}
\usepackage{longtable}
\usepackage{array}
\usepackage{hyperref}
\usepackage{setspace}

\setstretch{1.12}
\setlength{\parindent}{0pt}
\setlength{\parskip}{0.75em}
\emergencystretch=3em

\newcommand{\code}[1]{\texttt{#1}}
\newcommand{\metric}[1]{\texttt{#1}}
\newcommand{\yes}{\textsc{yes}}
\newcommand{\no}{\textsc{no}}

\title{GPT-5.5 Plan: Minimal-Endpoint Multi-Objective Cryomicroneedle Optimization}
\author{CryoMN ML Project}
\date{\today}

\begin{document}
\maketitle

\begin{abstract}
This document specifies the next-generation implementation plan for cryomicroneedle
(CryoMN) formulation optimization. The previous repository successfully used
literature transfer, Gaussian-process surrogate modeling, Bayesian optimization, and
wet-lab validation to improve mesenchymal stem-cell post-thaw viability. The next
project must add mechanical robustness without turning the experimental workflow into
an impractical mechanical-testing campaign. The proposed system therefore uses only
two optimization objectives, one required screening gate, one secondary mechanical
diagnostic, and one final qualitative insertion check. The main optimization objectives
are post-thaw viability and Instron-derived axial critical load. Intact patch formation
is modeled as a feasibility gate, initial stiffness is retained only as a diagnostic or
tie-breaker, and thumb insertion is recorded only qualitatively.
\end{abstract}

\section{Executive Summary}

The new project should be implemented as a cost-aware, minimal-endpoint,
multi-objective active-learning workflow. The purpose is not to discover the strongest
possible frozen material in isolation, but to find CryoMN formulations that preserve MSC
viability while forming intact, mechanically useful frozen microneedle patches.

The production objectives are:
\begin{itemize}
    \item maximize \metric{viability\_percent};
    \item maximize \metric{critical\_axial\_load\_N\_per\_needle}.
\end{itemize}

The required screening gate is:
\begin{itemize}
    \item \metric{intact\_patch\_formation\_pass}.
\end{itemize}

The only secondary mechanical endpoint is:
\begin{itemize}
    \item \metric{initial\_stiffness\_N\_per\_mm\_per\_needle}.
\end{itemize}

The final insertion check is:
\begin{itemize}
    \item \metric{thumb\_insertion\_pass}, recorded only as a qualitative yes/no
    result with optional notes or an image link.
\end{itemize}

The optimizer should not estimate insertion force, insertion margin, or Young's
modulus in version 1. Thumb insertion is deliberately uncontrolled and should not be
converted into a quantitative mechanical target.

\section{Endpoint Hierarchy}

\subsection{Screening Gate}

\metric{intact\_patch\_formation\_pass} is measured before mechanical testing and is
not a Pareto objective. It is a feasibility gate and classifier endpoint.

Default pass definition:
\begin{itemize}
    \item the CryoMN patch is visibly filled after molding;
    \item the patch demolds and handles as an intact frozen microneedle patch;
    \item there is no slurry-like collapse, gross bending, patch disintegration, or
    obvious unfrozen liquid phase;
    \item at least 90 of 100 tips are visually intact after demolding and handling.
\end{itemize}

If the laboratory later changes the mold geometry, the 90/100 rule should be stored as
a configurable threshold, not hard-coded into downstream modeling logic.

\subsection{Primary Objectives}

\begin{longtable}{@{}p{0.38\linewidth}p{0.54\linewidth}@{}}
\toprule
\textbf{Endpoint} & \textbf{Definition and use} \\
\midrule
\metric{viability\_percent} &
Post-thaw MSC viability, measured using the existing viability assay. This remains
the biological objective transferred from the current repository. \\
\metric{critical\_axial\_load\_N\_per\_needle} &
Maximum compressive load before buckling, fracture, sudden load drop, or force
plateau during Instron axial compression, divided by the number of needles actually
compressed. This is the mechanical objective. \\
\bottomrule
\end{longtable}

\subsection{Single Secondary Endpoint}

\metric{initial\_stiffness\_N\_per\_mm\_per\_needle} is calculated as the slope of
the early linear force-displacement region, divided by the number of needles
compressed. It should be stored, visualized, and optionally used as a diagnostic
covariate or tie-breaker. It is not a Pareto objective in version 1.

\subsection{Final Qualitative Check}

\metric{thumb\_insertion\_pass} is recorded only for finalists. It is a qualitative
yes/no result, optionally accompanied by free-text notes or an image link. Because the
applied thumb force is uncontrolled and unquantifiable, this endpoint must not be used
to calculate insertion force, insertion margin, or a continuous insertion objective.

\section{Instron 5942 Mechanical Protocol}

The Instron 5942 is a 0.5 kN single-column test system in the 5900 series. The relevant
capability for this project is controlled force-displacement testing in compression
using Bluehill-compatible load and extension acquisition. The CryoMN mechanical assay
should therefore be a low-temperature axial compression test against a rigid platen or
equivalent rigid contact surface.

The mechanical test should export raw force-displacement data for every specimen.
The importer will compute only:
\begin{itemize}
    \item \metric{critical\_axial\_load\_N\_per\_needle};
    \item \metric{initial\_stiffness\_N\_per\_mm\_per\_needle}.
\end{itemize}

The raw file and observation record must include the following metadata:
\begin{itemize}
    \item formulation identifier;
    \item experiment date;
    \item test temperature or cold-stage condition;
    \item number of needles compressed;
    \item crosshead speed;
    \item fixture/platen description;
    \item whether the curve had a clear peak, sudden drop, plateau, or no clear failure;
    \item operator notes.
\end{itemize}

Young's modulus should not be optimized in version 1. For frozen pyramidal CryoMN
arrays, apparent modulus would be confounded by patch base deformation, fixture
compliance, nonuniform tip contact, alignment, temperature drift, and thawing. The more
defensible version-1 measurements are peak axial load and early force-displacement
stiffness.

\section{Experimental Policy}

Each active-learning round should assume tight wet-lab capacity:
\begin{itemize}
    \item 12 viability screens;
    \item intact patch formation screening for all 12 if enough molded material exists;
    \item 3--4 Instron mechanical tests;
    \item thumb insertion only for finalists.
\end{itemize}

Mechanical testing should be selected from candidates with
\metric{intact\_patch\_formation\_pass = true}. One novelty or coverage exception per
round is allowed if it is scientifically useful, but the exception must be flagged in
the batch metadata so later analysis does not treat it as an ordinary mechanical
candidate.

\section{Constraints and Penalties}

\subsection{Ingredient Count}

Ingredient count is a soft penalty, not a hard cutoff. Candidates with more than 8
active ingredients may be exported if the acquisition value justifies them, but they
must be penalized and flagged.

For candidate $x$, let $n_{\mathrm{active}}(x)$ be the number of active ingredients
after applying the existing practical trace floor. The ingredient-count penalty should
be:
\[
P_{\mathrm{ingredients}}(x)
= \lambda_{\mathrm{ingredients}}\,
\max(0, n_{\mathrm{active}}(x)-8).
\]

The implementation should store both the raw active count and the penalty value.

\subsection{Single-Ingredient 500 mM Rule}

The 500 mM rule applies to each single molar ingredient independently. It is not a
total-solute-concentration penalty.

For each molar feature $x_j$, the concentration penalty should be:
\[
P_{\mathrm{500mM}}(x)
= \lambda_{\mathrm{500mM}}\,
\sum_{j \in \mathcal{M}}
\max(0, x_j - 0.5)^2,
\]
where $\mathcal{M}$ is the set of molar-valued ingredient features and concentrations
are represented in molar units. Percentage-valued ingredients should use separate
registry-defined caps.

The candidate generator may still emit a candidate with one ingredient above 0.5 M if
the acquisition value is compelling, but the row must be flagged as
\metric{single\_ingredient\_gt\_500mM = true}.

\subsection{Culture Media Exclusion}

Culture media, basal media, salts, and buffers remain excluded as formulation variables.
They may be recorded as fixed protocol context, but they must not appear as optimizable
features.

\section{Ingredient Registry}

The next implementation should introduce a central formulation registry. This registry
must define canonical names, feature names, units, synonyms, molecular weights or
conversion rules, optimization bounds, percentage caps, and whether the 500 mM penalty
applies.

The following ingredients must be added or activated:
\begin{longtable}{@{}p{0.28\linewidth}p{0.18\linewidth}p{0.45\linewidth}@{}}
\toprule
\textbf{Ingredient} & \textbf{Status} & \textbf{Implementation note} \\
\midrule
taurine & parsed but inactive & Activate as \metric{taurine\_M} if retained by registry rules. \\
sericin & parsed but inactive & Retain as percentage-valued \metric{sericin\_pct}. \\
raffinose & already present & Keep as \metric{raffinose\_M}. \\
methylcellulose & already present & Keep as \metric{methylcellulose\_pct}. \\
myo-inositol & new & Add as molar-valued \metric{myo\_inositol\_M}. \\
4-methoxyphenyl beta-D-glucopyranoside & new & Add as molar-valued \metric{methoxyphenyl\_glucopyranoside\_M}. \\
trehalose & already present & Keep as \metric{trehalose\_M}. \\
\bottomrule
\end{longtable}

The registry should be the single source of truth for feature generation. Parser,
candidate generator, penalty engine, batch exporter, and report code should all consume
the same registry.

\textbf{Current acquisition restrictions.} The legacy literature import should read
\code{data/processed/parsed\_formulations.csv}, but \metric{cooh\_pll},
\metric{ficoll}, and \metric{hes} are excluded from the active v2 feature set because
they cannot currently be acquired. Separately, ingredients that are valid scientific
candidates but temporarily unavailable in the lab should remain active registry
features and be excluded only at selection time through
\code{config\_v2/availability.yaml}. This temporary list currently contains
\metric{taurine\_M}, \metric{sericin\_pct}, \metric{raffinose\_M},
\metric{methylcellulose\_pct}, \metric{myo\_inositol\_M},
\metric{methoxyphenyl\_beta\_d\_glucopyranoside\_M}, and
\metric{trehalose\_M}.

\section{Data Transfer Policy}

The current repository transfers into the new project as viability evidence only. The
existing literature rows and 106 CryoMN wet-lab rows should seed the viability model,
but no mechanical labels should be inferred from legacy data.

Recommended source assumptions:
\begin{longtable}{@{}p{0.28\linewidth}p{0.34\linewidth}p{0.30\linewidth}@{}}
\toprule
\textbf{Source} & \textbf{Default noise} & \textbf{Use} \\
\midrule
Literature viability & 50 percentage points & Weak viability prior; 10$\times$ legacy wet-lab noise. \\
Legacy CryoMN wet-lab viability & 5 percentage points & Trusted viability correction signal. \\
New viability & 1 percentage point by default & Active biological observations; one fifth of legacy wet-lab noise unless overridden. \\
Intact patch formation & classifier likelihood & Feasibility model, not objective. \\
Instron critical load and stiffness & maximum of replicate SD, 15\% of mean, and instrument floor & Mechanical objective and diagnostic. \\
\bottomrule
\end{longtable}

\section{Algorithm}

\subsection{Cold Start}

Mechanical data start at zero, so the first phase should not rely immediately on a full
hypervolume acquisition function. Instead, use k-center or greedy diversity selection
among candidates that are predicted to be viable and likely to pass intact patch
formation. Continue this cold-start phase until at least 8 mechanical observations have
been collected.

Cold-start mechanical selection should:
\begin{itemize}
    \item prefer predicted intact-patch pass candidates;
    \item cover chemically distinct regions near promising viability predictions;
    \item include one optional novelty/coverage exception per round;
    \item avoid duplicate or near-duplicate formulations already tested mechanically.
\end{itemize}

\subsection{Main Acquisition}

After mechanical cold start, use noisy hypervolume-based multi-objective Bayesian
optimization. The production acquisition should be \code{qLogNEHVI}. The two objectives
used for hypervolume are normalized:
\begin{itemize}
    \item normalized \metric{viability\_percent};
    \item normalized \metric{critical\_axial\_load\_N\_per\_needle}.
\end{itemize}

\metric{intact\_patch\_formation\_pass} should enter as a feasibility classifier and
acquisition penalty. \metric{initial\_stiffness\_N\_per\_mm\_per\_needle} should not
enter the hypervolume objective set. It can be used only as a diagnostic, covariate, or
tie-breaker.

Evaluation baselines should include:
\begin{itemize}
    \item random or diversity sampling;
    \item \code{qLogNParEGO};
    \item known-noise or variance-aware \code{qLogNEHVI} variants.
\end{itemize}

\section{Computational Architecture}

\subsection{Directory Layout}

The next project should add new numbered workflow stages without breaking the legacy
viability pipeline:
\begin{itemize}
    \item \code{src/08\_multi\_objective/01\_build\_database/build\_database.py}:
    build the v2 database from legacy viability-only evidence;
    \item \code{src/08\_multi\_objective/02\_select\_candidates/select\_candidates.py}:
    score a candidate pool and export the next wet-lab slate;
    \item \code{src/08\_multi\_objective/03\_run\_round/run\_round.py}:
    generate pre/post round reviews, ingest filled wet-lab results, and generate the next slate;
    \item \code{src/08\_multi\_objective/helper/instron.py}:
    parse Bluehill/Instron CSV files into the round CSV;
    \item \code{src/08\_multi\_objective/helper/}: non-executable shared model,
    registry, instron, review, penalty, acquisition, and selection code.
\end{itemize}

\subsection{Decoupled Tables}

The new data layer should avoid forcing every endpoint into one wide validation CSV.
Use two persistent model-state tables and two selector-output CSV files:

\begin{longtable}{@{}p{0.25\linewidth}p{0.67\linewidth}@{}}
\toprule
\textbf{Table} & \textbf{Purpose} \\
\midrule
\code{formulations} &
One row per canonical formulation vector, including formulation ID, ingredient columns,
active ingredient count, 500 mM flags, percentage cap flags, and human-readable recipe. \\
\code{observations} &
One row per measured endpoint per formulation. Required columns include formulation ID,
endpoint name, value, uncertainty, replicate count, assay type, source, date, and notes. \\
\code{total\_candidate\_pool.csv} &
Full generated and scored candidate pool after temporary availability filtering. This
file is useful for audit and troubleshooting, but it is not wet-lab input and is not
persistent model state. \\
\code{next\_round\_candidates.csv} &
Per-round selector output containing candidate identity, formulation values,
predictions, feasibility diagnostics, mechanical-test recommendations, and blank
wet-lab result columns. This file is not persistent model state; stage 03 folds filled
results back into \code{observations}. \\
\bottomrule
\end{longtable}

Required observation endpoint names:
\begin{itemize}
    \item \metric{viability\_percent};
    \item \metric{intact\_patch\_formation\_pass};
    \item \metric{critical\_axial\_load\_N\_per\_needle};
    \item \metric{initial\_stiffness\_N\_per\_mm\_per\_needle};
    \item \metric{thumb\_insertion\_pass}.
\end{itemize}

\section{Acceptance Tests}

The implementation should include tests for the following scenarios:
\begin{itemize}
    \item intact-patch gate ingestion and classifier training;
    \item per-ingredient 500 mM penalties;
    \item ingredient-count soft penalties above 8 active ingredients;
    \item culture media and basal buffer exclusion;
    \item new ingredient registry entries and synonyms;
    \item Instron CSV parsing for peak load and initial stiffness;
    \item cold-start selection restricted to predicted intact-patch-pass candidates;
    \item \code{qLogNEHVI} candidate generation after seeded formation and mechanical data;
    \item LaTeX compilation of this planning document.
\end{itemize}

\section{Implementation Order}

\begin{enumerate}
    \item Add the registry and migrate parser/candidate code to consume it.
    \item Add decoupled formulation, observation, and batch table helpers.
    \item Implement the Instron importer and synthetic force-displacement test fixtures.
    \item Import legacy literature and CryoMN viability data as viability-only observations.
    \item Implement formation classifier training.
    \item Implement mechanical cold-start selection.
    \item Implement \code{qLogNEHVI} optimization once enough mechanical data exist.
    \item Add Pareto-front evaluation, normalized hypervolume, IGD, and diagnostic plots.
    \item Run regression tests to confirm the legacy single-objective viability pipeline remains usable.
\end{enumerate}

\section{Evidence Base}

The plan is grounded in four evidence streams. First, the current repository shows that
literature-only viability transfer is useful but insufficient without CryoMN-specific
wet-lab correction. Second, the Emerson et al. CPA-cocktail study supports noisy
hypervolume-based multi-objective Bayesian optimization over random search or purely
scalarized baselines for limited-experiment formulation discovery. Third, Instron
documentation confirms that the 5942/5900-series frame is appropriate for controlled
force-displacement testing in compression with Bluehill-compatible acquisition. Fourth,
microneedle mechanical-characterization literature supports extracting buckling or
fracture load and stiffness from load-displacement curves, while warning that speed,
alignment, base deformation, and number of compressed needles affect measured values.

\begin{thebibliography}{9}

\bibitem{instron5900}
Instron.
\textit{Out of Production 5900 Series Universal Testing Systems}.
\url{https://www.instron.com/en/products/testing-systems/universal-testing-systems/low-force-universal-testing-systems/5900-series/}

\bibitem{instron5940support}
Instron.
\textit{5940 Series Single Column Table Frames System Support}.
\url{https://www.instron.com/wp-content/uploads/2024/07/5940-single-column-table-frames-system-support.pdf}

\bibitem{mnmechanical}
Mechanical characterization of individual needles in microneedle arrays: factors
affecting compression test results.
\url{https://pdfs.semanticscholar.org/02d4/c6981289f07083ee62c04637c4ddd8e5c335.pdf}

\bibitem{polymermnfailure}
Park and Prausnitz.
\textit{Analysis of Mechanical Failure of Polymer Microneedles by Axial Force}.
\url{https://pmc.ncbi.nlm.nih.gov/articles/PMC3016089/}

\bibitem{emerson}
Emerson et al.
\textit{Accelerated Discovery of Cryoprotectant Cocktails via Multi-Objective Bayesian Optimization}.
Local repository PDF:
\code{docs/Accelerated Discovery of Cryoprotectant Cocktails via Multi-Objective Bayesian Optimization.pdf}.

\end{thebibliography}

\end{document}
